\documentstyle[%


righttag%


,11pt%


,amssymb%


,russian%               remove in Eng.


,izvru%                 Set izven% in Eng.


]{amsart}





%





\newcommand{\df}{\overset{\mathrm{def}}{=}}


\newcommand{\const}{\operatorname{const}}


\renewcommand{\le}{\leqslant}


\renewcommand{\ge}{\geqslant}





\renewcommand{\d}{\delta}


\renewcommand{\L}{{\cal L}}


\newcommand{\D}{D}


\renewcommand{\phi}{\varphi}


\newcommand{\A}{A}


\newcommand{\R}{{R}^1}


\newcommand{\SC}{C[a,b]}


\newcommand{\SL}{L[a,b]}


\newcommand{\tts}[1]{}


\renewcommand{\baselinestretch}{0.98}





\theoremstyle{plain}


\theorembodyfont{\it}


\newtheorem{thms}{\hskip\parindent Теорема}[section]


\newtheorem{lems}{\hskip\parindent Лемма}[section]





\theoremstyle{definition}


\theorembodyfont{\rm}


\newtheorem{notes}{\hskip\parindent Замечание}[section]


\newtheorem{defns}{\hskip\parindent Определение}[section]


\newtheorem{cors}{\hskip\parindent Следствие}[section]








%\hoffset=-15mm%                  For Eng.


\setcounter{page}{12}


\god{2001}


\nomer{6~(469)} %                        Remove in Eng.


%\nomer{Vol.\,45, No.\,6}%               Set in Eng.


%\str{\pageref{begin}--\pageref{end}}%  Set in Eng.


\udk{517.929}





\author{\it Е.И.\,Бравый}


% NB:   ^^ remove in Eng.


\title{О необходимых условиях


отрицательности\\ функции Грина двухточечной задачи}


\thanks{Работа выполнена при финансовой поддержке


Российского фонда фундаментальных исследований (96-15-96195, 99-01-01278) и


Конкурсного центра по исследованиям в области фундаментального


естествознания, Санкт-Петербург.}





\date{}


%\pagestyle{myheadings}%         Set in Eng.


\pagestyle{plain} %              Remove in Eng.


\begin{document}


%\thispagestyle{plain}%          Set in Eng.


\maketitle


\label{begin}





\section{Введение}





Вопрос о том, при каких условиях решение $x$ краевой


задачи


\begin{equation}


\label{00}


({\cal L} x)(t)=f(t),\ \; t\in[a,b];\quad x(a)=0,\ \; x(b)=0


\end{equation}


зависит от функции $f$ монотонно,


привлек внимание многих математиков и


рассматривался как для обыкновенных дифференциальных


операторов $\L$ (напр., \cite{Lev}, \cite{KSh}), так и для


функционально-дифференциальных.


Будем рассматривать линейный сингулярный


функционально-диф\-фе\-рен\-ци\-аль\-ный оператор $\L$


второго порядка. Если в этом случае задача (\ref{00}) имеет


единственное решение $x=Gf$ при любой правой части $f$, то


линейный оператор $G$ называется оператором Грина краевой задачи


(\ref{00}). Вопрос о монотонной зависимости решения задачи


(\ref{00}) от правой части сводится к условиям монотонности


оператора Грина $G$ краевой задачи.





Линейный оператор $A:X\to Y$, где $X$ и $Y$ ---


полуупорядоченные пространства, будем называть {\it


изотонным} ({\it антитонным}) и обозначать $A\ge0$ ($A\le0$), если


$Ax_2\ge Ax_1$ ($Ax_2\le Ax_1$) при всех $x_1\le x_2$,


$x_1,x_2\in X$.





Введем следующие обозначения:


$\R$ --- пространство вещественных чисел;


$C[a,b]$ --- пространство непрерывных функций $x:[a,b]\to


\R$ с нормой $\|x\|_C=\max\limits_{t\in[a,b]}|x(t)|;$


$L[a,b]$ --- пространство классов эквивалентности


суммируемых функций $z:[a,b]\to


\R$ с нормой $\|z\|_L=\int\limits_a^b|z(t)|\,dt;$


$\chi_e$ --- характеристическая функция множества $e\subset\R$;


$\rho(K)$ --- спектральный радиус оператора $K$;


$I$ --- тождественный оператор в соответствующем пространстве.





B пространствах $C[a,b]$ и $L[a,b]$ будем использовать


естественную полуупорядоченность.





Для линейных обыкновенных дифференциальных уравнений второго


порядка антитонность оператора Грина задачи (\ref{00})


эквивалентна наличию положительного решения


однородного уравнения $\L x=0$ (напр., \cite{H}, \cite{AD0}).


В этом случае оператор Грина


будет строго антитонен (см. определение \ref{def2} далее). Необходимые


условия антитонности оператора Грина задачи (\ref{00}) для


функционально-дифференциальных уравнений


изучены сравнительно мало:


если оператор $\L$ имеет вид


$$


\L x\df\pi \Ddot x - Tx,


$$


где $T: C[a,b]\to L[a,b]$, $T\le0$ и


$\pi(t)=(t-a)^{\mu_1}(b-t)^{\mu_2}$, $0\le\mu_1,\mu_2\le1$,


то необходимым и достаточным условием антитонности


оператора Грина $G$ является неравенство


$\rho(G_0T_{C[a,b]\to C[a,b]})<1$ \cite{AMR}, \cite{Lab}. Здесь $G_0$


--- оператор Грина задачи


\begin{equation}


\label{0}


(\L_0x)(t)\df\pi(t) \Ddot x(t)=f(t),\ \; t\in[a,b];\quad


x(a)=0,\ \; x(b)=0.


\end{equation}


Оператор Грина $G_0:L[a,b]\to C[a,b]$ антитонен и имеет представление


$$


(G_0z)(t)\df \int_a^bG_0(t,s)z(s)\,ds\equiv 


\int_a^t \frac{(t-b)(s-a)}{(b-a)\pi(s)}z(s)\,ds+


\int_t^b \frac{(s-b)(t-a)}{(b-a)\pi(s)}z(s)\,ds,\quad t\in[a,b].


$$





В случае, когда оператор $\L$ имеет вид


$$


\L x\df\pi \Ddot x + T^+x - T^-x,


$$


где $T^+,T^-: C[a,b]\to L[a,b]$, $T^+,T^-\ge0$,


для получения достаточных условий антитонности оператора


Грина требуется условие антитонности оператора Грина $G_1$


задачи


$$


(\L_1 x)(t)\df\pi(t) \Ddot x(t)-(T^-x)(t)=f(t),\ \; t\in[a,b];\quad


x(a)=0,\ \; x(b)=0.


$$





Эта задача, конечно, имеет и


самостоятельный интерес. В случае, когда функция $\pi(t)$ обращается в


нуль на отрезке $[a,b]$, задачу будем называть сингулярной.


Достаточные условия антитонности оператора Грина $G_1$ получены


в \cite{Kob}--\cite{AD2} для операторов без сингулярностей


и в \cite{AAB} --- для операторов с сингулярностями.


Утверждение о необходимом условии антитонности оператора Грина


$G_1$ при $\pi(t)\equiv1$ и вольтерровом операторе $T$ приведено


в \cite{Lab2}.





Во всех этих работах условия антитонности оператора


$G_1$ получены в тех случаях, когда вронскиан фундаментальной


системы решений однородного уравнения $\L_1 x=0$ не обращается в


нуль на интервале $(a,b)$. Тогда


функционально-дифференциальное уравнение


эквивалентно некоторому обыкновенному дифференциальному


уравнению \cite{AD}.





Следующий пример краевой задачи показывает, что необращение


в нуль вронскиана фундаментальной


системы решений уравнения $\L_1 x=0$ вовсе не является


необходимым для антитонности оператора $G_1$.





Рассмотрим краевую задачу


\begin{equation}


\label{1}


(\L_1x)(t)\df\Ddot


x(t)-p(t)x(h(t))=f(t),\ \; t\in[-1,1];\quad


x(-1)=0,\ \; x(1)=0,


\end{equation}


где


\begin{align*}


p(t)&=


\begin{cases}


2/{c^2_1}& \text{ при }t\in[-1,0);\\


{2}/{c^2_2}& \text{ при }t\in[0,1],


\end{cases}


\quad c_1\in[-1,0),\quad c_2\in(0,1],\\


%


h(t)&=


\begin{cases}


{c_1}& \text{ при }t\in[-1,0);\\


c_2& \text{ при }t\in[0,1].


\end{cases}


\end{align*}


Множество решений однородного уравнения $\L_1 x=0$


двумерно, функции $x_1(t)=t^2\chi_{[0,1]}(t)$,


$x_2(t)=t^2\chi_{[-1,0]}(t)$ образуют фундаментальную систему


решений. Вронскиан этой системы решений --- тождественный нуль


на всем отрезке $[-1,1]$. Таким образом, в частности, ни одна из


задач $\L_1x=f$, $x(\tau)=0$, $\Dot x(\tau)=0$ при всех


$\tau\in[-1,1]$ не является однозначно разрешимой. Тем не менее


задача (\ref{1}) однозначно разрешима. Непосредственным


построением оператора Грина доказывается





\begin{thms}


Оператор Грина задачи $(\ref{1})$ антитонен


тогда и только тогда, когда $-1\le c_1\le-1/2$


и $1/2\le c_2\le1$.


\end{thms}





Нас интересуют необходимые условия антитонности оператора


Грина $G$ задачи


\begin{equation}


\label{11}


(\L x)(t)\df\pi(t) \Ddot x(t)-(Tx)(t)=f(t),\ \; t\in[a,b];\quad


x(a)=0,\ \; x(b)=0,


\end{equation}


где $T:C[a,b]\to L[a,b]$ --- линейный ограниченный оператор,


$\pi(t)=(t-a)^{\mu_1}(b-t)^{\mu_2}$,


$0\le\mu_1, \mu_2\le1$.


Решение задачи (\ref{11}) ищется в банаховом пространстве $\D$ всех функций


$x:[a,b]\to \R$, обладающих следующими свойствами:





1) функция $x$ абсолютно непрерывна на $[a,b]$;





2) функция $\Dot x$ абсолютно непрерывна на каждом


замкнутом промежутке из $(a,b)$;





3) $\int\limits_a^b\pi(t)|\Ddot x(t)|\,dt<\infty$.





Норма в пространстве $\D$ определена равенством


$$


\|x\|_{\D}=|x(a)|+|x(b)|+\int_a^b\pi(t)|\Ddot x(t)|\,dt.


$$


В работе \cite{Lab} показано, что пространство $\D$ непрерывно вложено в


пространство абсолютно непрерывных функций и,


следовательно, в пространство $C[a,b]$. Кроме того, там же


доказано, что задача $\pi\Ddot x=f$, $x(a)=c_1$,


$x(b)=c_2$ имеет единственное решение $x\in\D$ при всех


$f\in L[a,b]$, $c_1, c_2\in R^1$. Если $\pi(t)\equiv1$, то


пространство $\D$ совпадает с пространством всех функций с абсолютно


непрерывной производной.





Оператор Грина $G$ как линейный ограниченный оператор, действующий из


пространства суммируемых функций в пространство непрерывных


функций, имеет интегральное представление. Ядро этого


представления, функция $G(t,s)$, $t,s\in[a,b]$,


называется функцией Грина краевой задачи. Очевидно, что при


каждом $t\in[a,b]$ функция $G(t,\cdot)$ может быть произвольно


изменена на любом множестве меры нуль. Таким образом, любое


утверждение относительно функции Грина следует понимать так:


среди эквивалентных (т.\,е. задающих один и тот же оператор)


функций Грина найдется функция Грина, удовлетворяющая данному


утверждению. Очевидно, что антитонность оператора Грина


эквивалентна неположительности функции Грина.





\begin{defns}


Следуя \cite{AAB}, будем говорить, что оператор $\L$ обладает


свойством $\A$, если краевая задача


$$


\L x=f,\quad x(\tau)=0,\quad \Dot x(\tau)=0


$$


однозначно разрешима при всех $f\in L[a,b]$ и ее оператор Грина


изотонен при всех $\tau\in(a,b)$.


\end{defns}





\begin{notes}


Если оператор $\L$ обладает свойством $\A$, то фундаментальная


система решений однородного уравнения $\L x=0$


двумерна и ее вронскиан отличен от нуля в каждой


точке интервала $(a,b)$.


\end{notes}





\begin{defns}\label{def2}


Назовем антитонный оператор Грина $G$ задачи (\ref{11})


строго антитонным, если для каждой неотрицательной не


эквивалентной нулю суммируемой функции $z$ выполнено


неравенство


$$


(Gz)(t)<0 \quad \text{при всех}\ \;t\in(a,b).


$$


\end{defns}





\begin{notes}


Оператор Грина задачи (\ref{11}) строго антитонен тогда и только


тогда, когда для этой задачи справедлив ``принцип


максимума'': при неотрицательной не равной тождественно нулю


правой части решение краевой задачи достигает своего максимума


только на концах отрезка.


\end{notes}





В следующем разделе получим необходимые условия антитонности


оператора Грина. В \S\,3 исследуется связь между


антитонностью оператора Грина и свойством $\A$. В \S\,4


при некоторых дополнительных условиях доказана


эквивалентность строгой антитонности


оператора Грина и свойствa $\A$.





\section{Необходимые условия антитонности оператора Грина}





Определим при $s\in(a,b)$ непрерывные функции $d^a_s$ и $d^b_s$ равенствами


\begin{alignat*}{2}


d^a_s(t)&=\chi_{[a,s]}(t)\dfrac{s-t}{s-a},&\quad t&\in[a,b],\\


d^b_s(t)&=\chi_{[s,b]}(t)\dfrac{t-s}{b-s},&\quad t&\in[a,b].


\end{alignat*}





\begin{thms}\label{t2}


Пусть выполнены условия


\begin{equation}\label{9}


\lim_{s\to b-}


\int_a^b\left|(Td^b_s)(\theta)\right|\,d\theta=0,


\quad


\lim_{s\to a+}


\int_a^b\left|(Td^a_s)(\theta)\right|\,d\theta=0.


\end{equation}


Тогда, если краевая задача $(\ref{11})$ имеет единственное решение


$x\in\D$ при всех $f\in L[a,b]$ и ее


оператор Грина антитонен, то решения задач


\begin{equation}\label{7}


(\L x)(t)=0,\ \; t\in[a,b];\quad


x(a)=0,\ \; x(b)=1


\end{equation}


и


\begin{equation}\label{8}


(\L x)(t)=0,\ \; t\in[a,b];\quad


x(a)=1,\ \; x(b)=0


\end{equation}


неотрицательны.


\end{thms}





\begin{notes}


Если $T$ --- антитонный оператор, то решения задач (\ref{7}) и


(\ref{8}) положительны в интервале $(a,b)$ и без условий (\ref{9}).


\end{notes}





\begin{notes}


В (\cite{KVP}, с.\,317) доказано, что любой регулярный линейный


ограниченный оператор $T: C[a,b]\to L[a,b]$, т.\,е. представимый


в виде разности двух изотонных операторов, имеет представление в


виде интеграла Стилтьеса


\begin{equation}\label{st}


(Tx)(t)=\int_a^bx(s)\,d_sr(t,s),


\end{equation}


где $\int\limits_a^b|r(t,s)|\,dt<\infty$ при почти всех $s\in[a,b]$,


функции $r(t,\cdot)$ имеют ограниченное изменение на $[a,b]$ и


$\int_a^b \underset{s\in[a,b]}{\mathrm{Var}}{} r(t,s)\, dt<\infty$.


При этом, если оператор $T$ изотонен, то при почти всех $t\in[a,b]$ функции


$r(t,\cdot)$ не убывают на $[a,b]$.





Если оператор $T: C[a,b]\to L[a,b]$ имеет представление (\ref{st}),


то условия (\ref{9}) выполнены тогда и только тогда, когда для


почти всех $t\in[a,b]$ выполнены равенства


\begin{equation}\label{n9}


\lim_{s\to a+}r(t,s)=r(t,a),\quad \lim_{s\to b-}r(t,s)=r(t,b).


\end{equation}


Это условие использовалось и в работе \cite{Lab2}.


\end{notes}





\begin{pf}


Задача (\ref{11}) эквивалентна уравнению в пространстве $C[a,b]$


\begin{equation}


\label{22}x=G_0Tx+G_0f,


\end{equation}


где, как и ранее, $G_0$ --- оператор Грина задачи (\ref{0}).


Из того, что уравнение (\ref{22}) имеет единственное решение


при всех $f\in L$, следует, что оператор $I-G_0T:C[a,b]\to C[a,b]$


обратим \cite{MR}. Обратный оператор $B\df(I-G_0T)^{-1}$


связан с оператором Грина $G$ равенством $Gf=BG_0f$ для каждой


$f\in L[a,b]$.





Операторы $G_0$ и $G$ антитонные. Кроме того, оператор $G_0$ отображает


множество неотрицательных функций на все


множество функций из $\D$, равных нулю в точках $a$ и $b$ и


выпуклых вниз. Таким образом, оператор $B$ отображает выпуклые


вверх неотрицательные функции с абсолютно непрерывной


производной, равные нулю в точках $a$ и $b$, в неотрицательные


функции. Оператор $B$ действует непрерывно в пространстве


$C[a,b]$. Поэтому он отображает любую выпуклую вверх


непрерывную неотрицательную функцию, обращающуюся в нуль в


точках $a$ и $b$, в неотрицательную функцию.





Для каждого $s\in(a,b)$ обозначим


\begin{gather*}


f_s(t)=


\dfrac{t-a}{b-a}-\chi_{[s,b]}(t)\dfrac{t-s}{b-s},\quad t\in[a,b], \\


%


g_s(t)=


\dfrac{b-t}{b-a}-\chi_{[a,s]}(t)\dfrac{s-t}{s-a},\quad t\in[a,b].


\end{gather*}


Кроме того, введем обозначения


$$


f_b(t)=\dfrac{t-a}{b-a},\quad


g_a(t)=\dfrac{b-t}{b-a},\quad t\in[a,b].


$$





Уравнение $x=G_0Tx+f$ имеет единственное решение $x\in C[a,b]$ при


любом $f\in C[a,b]$. Поэтому единственным образом


определены функции $x_b$, $x_s$, $y_a$, $y_s$, удовлетворяющие


соответственно равенствам


\begin{gather*}


x_b=G_0Tx_b+f_b,\quad x_s=G_0Tx_s+f_s,\\


y_a=G_0Ty_a+g_a,\quad y_s=G_0Ty_s+g_s.


\end{gather*}





Так как оператор $(I-G_0T)^{-1}$ отображает любую


непрерывную выпуклую вверх функцию, равную нулю в точках $a$ и


$b$, в неотрицательную функцию, то $x_s\ge0$, $y_s\ge0$


при всех $s\in(a,b)$.





Из того, что задача (\ref{11}) однозначно разрешима, следует, что


обратим оператор $I-TG_0:L[a,b]\to L[a,b]$. Поэтому решение


уравнения $x=G_0Tx+f$ можно записать в виде


$$


x=G_0(I-TG_0)^{-1}Tf+f


$$


(действительно, $Tx=TG_0Tx+Tf$, следовательно, $Tx=(I-TG_0)^{-1}Tf$).





Таким образом,


$x_b-x_s=G_0T(x_b-x_s)+f_b-f_s$.


Поэтому


$x_b-x_s=G_0(I-TG_0)^{-1}T(f_b-f_s)+f_b-f_s$.


Аналогично получаем


$y_a-y_s=G_0(I-TG_0)^{-1}T(g_a-g_s)+g_a-g_s$.


Отметим, что $d_s^a=g_a-g_s$, $d_s^b=f_b-f_s$. Поэтому


условия (\ref{9}) дают следующие равенства:


$y_a-g_a=\lim\limits_{s\to a+}(y_s-g_s)$, $x_b-f_b=\lim\limits_{s\to


b-}(x_s-f_s)$


или


$y_a=\lim\limits_{s\to a+}(y_s+g_a-g_s)$,


$x_b=\lim\limits_{s\to b-}(x_s+f_b-f_s)$


(здесь пределы рассматриваются в пространстве $C[a,b]$).





Так как $g_a\ge g_s$, $f_b\ge f_s$ и $y_s\ge0$, $x_s\ge0$, то


$y_a\ge0$, $x_b\ge0$.


Более того, $y_a(t)=\lim\limits_{s\to a+}y_s(t)$ для каждого $t\in(a,b]$;


$x_b(t)=\lim\limits_{s\to b-} x_s(t)$ для каждого $t\in[a,b)$.


Очевидно, $y_a$ удовлетворяет задаче


(\ref{8}) и функция $x_b$ удовлетворяет задаче (\ref{7}).


\end{pf}





\begin{notes}


В условиях теоремы \ref{t2} выполнены


равенства


\begin{equation}\label{zv}


y_a(t)=-\lim\limits_{s\to


a+}\dfrac{G(t,s)\pi(s)}{s-a},\quad x_b(t)=-\lim\limits_{s\to


b-}\dfrac{G(t,s)\pi(s)}{b-s},\quad t\in[a,b],


\end{equation}


где $y_a$, $x_b$ --- решения задач (\ref{8}), (\ref{7}) соответственно.





Действительно, имеют место соотношения


$$


y_s(t)=-\dfrac{G(t,s)\pi(s)}{s-a},\quad x_s(t)=-\dfrac{G(t,s)\pi(s)}{b-s},


\quad t\in[a,b]


$$


(прямой подстановкой можно убедиться, что, напр., функция


$$


x(t)=-\int_a^b y_s(t)(s-a)/\pi(s) f(s)\,ds,\quad t\in[a,b],


$$


удовлетворяет однородным краевым условиям и уравнению


$\pi\Ddot x=Tx+f$ при ограниченном операторе $T: C[a,b]\to


L[a,b]$). Отсюда следуют равенства $(\ref{zv})$.


\end{notes}





\begin{cors}\label{sl1}


Оператор Грина задачи


\begin{equation}\label{zsl1}


(\L_1x)(t)\df


\pi(t)\Ddot x(t)-\displaystyle\int_a^b R(t,s)x(s)\,ds=f(t),


\ \; t\in[a,b];\quad


x(a)=0,\ \; x(b)=0


\end{equation}


не является антитонным, если


$R(t,s)\ge\varepsilon$, $t,s\in[a,b]$, при некотором


$\varepsilon>0$, и либо $\mu_1=1$, либо $\mu_2=1$.


\end{cors}





\begin{pf}


Предположим, что $\mu_1=1$. Покажем, что решение $x_b$ задачи


$\L_1x=0$, $x(a)=0$, $x(b)=1$ не может быть


неотрицательным в $(a,b)$. Если $x_b(t)\ge0$ при


$t\in(a,b)$, то найдется такое число $\d>0$, что $\pi(t)\Ddot


x_b(t)\ge\d$ при почти всех $t\in[a,b]$. Поэтому $x_b(t)\le


y(t)$, $t\in[a,b]$, где $y$ --- решение задачи $(t-a)\Ddot


y(t)=\d_1\df \d (b-a)^{-\mu_2}$, $t\in[a,b]$; $y(a)=0$,


$y(b)=1$. Имеем


$$


y(t)=\d_1(t-a)\ln\frac{t-a}{b-a}+\frac{t-a}{b-a},\quad


t\in(a,b].


$$


Так как $y(a)=0$ и $\lim\limits_{t\to a+}\Dot y(t)=-\infty$, то


функция $y$, а следовательно, и функция $x_b$ принимают


отрицательные значения в $(a,b)$. Таким образом, по теореме


\ref{t2} оператор Грина задачи (\ref{zsl1}) не может быть


антитонным.





В случае $\mu_2=1$ доказательство аналогично.


\end{pf}





\section{Свойство $A$ и антитонность оператора Грина}





В работе \cite{AAB} показано, что если оператор $T$ изотонен и


оператор $\L$ обладает свойством $\A$, то оператор Грина задачи


(\ref{11}) является строго антитонным. Там же показано, что


оператор $\L$ при изотонном операторе $T$ обладает свойством


$\A$ тогда и только тогда, когда решения задач (\ref{7}), (\ref{8})


положительны в интервале $(a,b)$.





Обратить теорему \ref{t2} нельзя: пример краевой 


задачи (\ref{1}) показывает, что условия (\ref{9}) вместе с


однозначной разрешимостью двухточечной задачи и


неотрицательностью решений задач (\ref{7}), (\ref{8}) еще не


гарантируют антитонность оператора Грина задачи (\ref{11}) даже


при изотонном операторе $T$ в случае, когда $-1/2<c_1<0$ или $0\le


c_2<1/2$.





\begin{thms}


Пусть оператор $T$


изотонен, выполнены условия $(\ref{9})$, фундаментальная


система решений уравнения $\L x=0$ двумерна и ее вронскиан не


обращается в нуль на интервале $(a,b)$. Тогда оператор Грина


задачи $(\ref{11})$ антитонен тогда и только тогда, когда


выполнено свойство $\A$.


\end{thms}





\begin{pf}


В \cite{AAB} показано,


что при изотонном операторе $T$ из свойства $\A$ следует


строгая антитонность оператора Грина.





Если вронскиан фундаментальной системы решений


уравнения $\L x=0$ не обращается в нуль на интервале $(a,b)$,


то при изотонном операторе $T$ неотрицательные решения задач


(\ref{7}), (\ref{8}) положительны на интервале $(a,b)$. Таким


образом, как показано в \cite{AAB}, оператор $\L$


обладает свойством $\A$.


\end{pf}





\begin{thms}


Пусть оператор $T$ изотонен,


выполнены условия $(\ref{9})$ и оператор Грина задачи


$(\ref{11})$ антитонен. Тогда, если


вронскиан фундаментальной системы решений уравнения $\L x=0$


отличен от нуля в некоторой точке $\tau\in(a,b)$, то оператор Грина


задачи


\begin{equation}\label{tbvp}


(\L x)(t)=f,\ \; t\in[a,b];\quad


x(\tau)=0,\ \; \Dot x(\tau)=0


\end{equation}


изотонен.


\end{thms}





\begin{pf}


Краевая задача (\ref{tbvp}) эквивалентна уравнению


$$


x(t)=(W_\tau Tx)(t)+(W_\tau f)(t),\quad


t\in[a,b],


$$


в пространстве $C[a,b]$. Здесь $W_\tau:L[a,b]\to C[a,b]$ ---


оператор Грина задачи


$$


(\L_0x)(t)\df\pi(t) \Ddot x(t)=f(t),\ \;


t\in[a,b];\quad x(\tau)=0,\ \; \Dot x(\tau)=0.


$$


Отметим, что оператор $W_\tau$ изотонен.





Решения задач (\ref{7}), (\ref{8}) неотрицательны


(теорема \ref{t2}). Оператор $T$ изотонен, следовательно, решение первой


задачи не убывает, а решение второй из этих задач не возрастает,


и, кроме того, по условию теоремы оба эти решения не обращаются в


нуль в точке $\tau$. Поэтому существует такая линейная комбинация $u$


этих решений с неотрицательными коэффициентами, что


$u(t)>0$ при всех $t\in[a,b]$ и $\Dot u(\tau)=0$.


Следовательно, положительная функция $u$ удовлетворяет равенству


$$


u(t)=(W_\tau Tu)(t)+u(\tau),\quad t\in[a,b],


$$


причем $u(\tau)>0$. Таким образом, спектральный радиус изотонного


оператора $W_\tau T$ меньше единицы, и оператор Грина задачи (\ref{tbvp})


$$


W=(I-W_\tau T)^{-1}W_\tau=(I+W_\tau T+(W_\tau


T)^2+\dotsb)W_\tau


$$


изотонен (напр., \cite{AR}).


\end{pf}





\section{Условия строгой антитонности оператора Грина}





Как уже было отмечено, в \cite{AAB} показано, что положительность


решений задач


(\ref{7}) и (\ref{8}) достаточна для строгой антитонности оператора


Грина задачи (\ref{11}) при изотонном $T$. Покажем, что при


некоторых дополнительных предположениях эти условия также и


необходимы для строгой антитонности оператора Грина,


что означает эквивалентность свойства $\A$ и строгой антитонности оператора


Грина.





\begin{thms}\label{t5}


Пусть оператор $T:C[a,b]\to L[a,b]$


имеет представление $(\ref{st})$, в котором при


почти всех $t\in(a,b)$


\begin{align*}


r(t,s)&=r(t,a)\text{ при всех $s$ из некоторой


правой окрестности точки $a;$}\\


%


r(t,s)&=r(t,b)\text{ при всех $s$ из некоторой


левой окрестности точки $b$}.


\end{align*}


Тогда, если оператор


Грина $G$ задачи $(\ref{11})$


строго антитонен, то решения задач $(\ref{7})$ и $(\ref{8})$ положительны


в интервале $(a,b)$.


\end{thms}





\begin{notes}


В условиях теоремы \ref{t5} не требуется


изотонность оператора $T$ и выполнены условия (\ref{n9}).


\end{notes}





Для доказательства теоремы потребуются вспомогательные утверждения.





\begin{lems}


\label{l1}


Для каждого $t\in[a,b]$ функция Грина $G(t,\cdot)$ краевой задачи


$(\ref{11})$ удовлетворяет


при почти всех $s\in[a,b]$ уравнению


\begin{equation}\label{g1}


G(t,s)-(TG_0)^*(G(t,\cdot))(s)=G_0(t,s).


\end{equation}


\end{lems}





\begin{pf}


Краевая задача $\pi\Ddot x-Tx=f$, $x(a)=0$, $x(b)=0$ может быть с


помощью представления $x=G_0(\pi\Ddot x)$ записана в виде $Q(\pi\Ddot


x)=f$, где $Q=I-TG_0$, $Q:L[a,b]\to L[a,b]$.





Тогда имеем


$$


G=G_0(I-TG_0)^{-1}


$$


и


$$


GQf=G_0Q^{-1}Qf=G_0f,\quad f\in L[a,b].


$$





При каждом $t\in[a,b]$ и каждой функции $f\in L[a,b]$ имеет


место равенство


$$


\int_a^bG(t,s)(Qf)(s)\,ds=\int_a^bQ^*[G(t,\cdot)](s)f(s)\,ds,


$$


где сопряженный оператор $Q^*$ действует в пространстве


ограниченных в существенном функций.








Таким образом, при каждом $t\in[a,b]$ справедливо


равенство


$$


\int_a^b[G(t,s)-(TG_0)^*(G(t,\cdot))(s)]f(s)\,ds


=\int_a^bG_0(t,s)f(s)\,ds


$$


при всех $f\in L[a,b]$. Отсюда следует равенство


(\ref{g1}).


\end{pf}





С помощью прямых вычислений получаем следующее утверждение.





\begin{lems}


\label{l2}


Пусть оператор


$T: C[a,b]\to L[a,b]$ имеет представление $(\ref{st})$. Тогда для каждой


функции $f\in L[a,b]$ при почти всех $s\in[a,b]$ выполнено


равенство


\begin{equation*}%\label{g2}


(TG_0)^*(f)(s)=-\dfrac{1}{\pi(s)}


\int_a^bK(s,\xi)f(\xi)\,d\xi,


\end{equation*}


где


\begin{equation}\label{K}


K(s,\xi)=\int_a^b(r(\xi,\theta)-r(\xi,a))\,d\theta


\frac{s-a}{b-a}-\int_a^s(r(\xi,\theta)-r(\xi,a))\,


d\theta.


\end{equation}


\end{lems}





\begin{pf*}{Доказательство теоремы \ref{t5}}





При фиксированном $t\in(a,b)$ из лемм \ref{l1} и \ref{l2}


следует


\begin{equation}\label{dzv}


\pi(s)G(t,s)=\phi(t,s)-\psi(t,s),


\end{equation}


где


\begin{align*}


\phi(s)&=


\begin{cases}


\dfrac{(t-b)(s-a)}{b-a}& \text{


при }a\le s\le t<b;\\


\dfrac{(t-a)(s-b)}{b-a}& \text{ при }


b\ge s>t\ge a,


\end{cases}


\\


\psi(t,s)&=


\int_a^bK(s,\xi)G(t,\xi)\,d\xi,\quad s\in[a,b],


\end{align*}


функция $K(s,\xi)$ определена равенством (\ref{K}).





Поскольку $r(\xi,\cdot)=\const$ при почти всех $\xi$ в


некоторых окрестностях точек $a$ и $b$, то в этих


окрестностях и при этих $\xi$


выполняется равенство


$$


\dfrac{\partial K(s,\xi)}{\partial


s}=\int_a^b(r(\xi,\theta)-r(\xi,a))\,d\theta\frac{1}{b-a}.


$$


Таким образом, функция $\psi(t,\cdot)$, а следовательно,


и функция $G(t,\cdot)\pi(\cdot)$ линейны в некоторых


окрестностях точек $a$ и $b$ при каждом $t\in(a,b)$. Так как


оператор $G$ строго антитонный, то в некоторых окрестностях


точек $a$ и $b$ линейная функция $G(t,\cdot)\pi(\cdot)$ не равна


тождественно нулю. Из равенств (\ref{zv}) заключаем, что решения


задач (\ref{7}) и (\ref{8}) положительны в $(a,b)$.


\end{pf*}





\begin{cors}


Оператор Грина задачи


$$


(\L_1x)(t)\df


\pi(t)\Ddot x(t)-p(t)x(h(t))=f(t),\ \; t\in[a,b];\quad


x(a)=0,\ \; x(b)=0


$$


не является строго антитонным, если $p(t)\ge\varepsilon$,


$h(t)\in[a+\varepsilon,b-\varepsilon]$, $t\in[a,b]$, при некотором


$\varepsilon>0$, и либо $\mu_1=1$, либо $\mu_2=1$.


\end{cors}





Доказательство аналогично доказательству следствия \ref{sl1}.








Существенность условий теоремы \ref{t5} показывает следующий


пример краевой задачи, обобщающей задачу (\ref{1}):


\begin{equation}\label{n1}


\begin{array}{c}


(\L x)(t)\df\Ddot x(t)-2


\begin{cases}


\frac{\int\limits_{-1}^0v(\tau)x(\tau)\,d\tau}


{\int\limits_{-1}^0v(\tau)\tau^2\,d\tau}, & t\in[-1,0],\\


\frac{\int\limits_{0}^1v(\tau)x(\tau)\,d\tau}


{\int\limits_{0}^1v(\tau)\tau^2\,d\tau}, & t\in(0,1],


\end{cases}


=f(t),\ \; t\in[-1,1];\quad


x(-1)=0,\ \; x(1)=0,


\end{array}


\end{equation}


где неотрицательная функция $v$ суммируема.


Здесь соответствующий оператор $T$ имеет представление (\ref{st}):


$$


R(t,s)=


\begin{cases}


2\frac{\int\limits_{-1}^sv(\tau)x(\tau)\,d\tau}


{\int\limits_{-1}^0v(\tau)\tau^2\,d\tau},&\ t\in[-1,0],\ s\in[-1,0];\\


%


2\frac{\int\limits_{-1}^0v(\tau)x(\tau)\,d\tau}


{\int\limits_{-1}^0v(\tau)\tau^2\,d\tau},&\ t\in[-1,0],\ s\in(0,1];\\


0,&\ t\in(0,1],\phantom{-}\ s\in[-1,0];\\


2\frac{\int\limits_{0}^sv(\tau)x(\tau)\,d\tau}


{\int\limits_{0}^1v(\tau)\tau^2\,d\tau},&\ t\in(0,1],\phantom{-}\ s\in(0,1].


\end{cases}


$$





\begin{thms}


Пусть функция $v$ положительна в некоторой правой окрестности точки


$t=-1$ и в некоторой левой окрестности точки


$t=1$, и выполнены неравенства


$$


\int_{-1}^0(2\tau+1)v(\tau)\,d\tau\le0,\quad


\int_0^1(2\tau-1)v(\tau)\,d\tau\ge0.


$$


Тогда оператор Грина задачи $(\ref{n1})$ строго антитонен.


\end{thms}





Утверждение доказывается с помощью прямого построения функции


Грина. Условия


теоремы \ref{t5} здесь не выполнены. Фундаментальная


система решений однородного уравнения $\L x=0$ совпадает с


фундаментальной системой решений уравнения в краевой задаче


(\ref{1}), таким образом, вронскиан фундаментальной системы


тождественно равен нулю и решения задач (\ref{7}), (\ref{8}) не


являются положительными на $(-1,1)$. Тем не менее оператор Грина


может быть строго антитонным.


\medskip





Теперь получим необходимые условия для антитонного, но не


строго антитонного оператора Грина $G$.





\begin{thms}\label{t7}


Пусть для линейного изотонного оператора


$T:C[a,b]\to L[a,b]$ выполнены условия $(\ref{n9})$. Предположим,


что оператор Грина $G$ задачи $(\ref{11})$ антитонен, но не


является строго антитонным. Тогда вронскиан $w$ фундаментальной


системы решений однородного уравнения $\L x=0$


имеет нули на интервале $(a,b)$,


причем, если $w(t_0)=0$ в некоторой точке $t_0\in(a,b)$, то либо $w(t)=0$


при $t\in[a,t_0]$, либо $w(t)=0$ при $t\in[t_0,b]$.


\end{thms}





Для доказательства потребуются следующие утверждения, проверяемые


непосредственными вычислениями.





\begin{lems}


\label{l3}


Пусть оператор $T:C[a,b]\to L[a,b]$


имеет представление $(\ref{st})$ и изотонен. Тогда функция


$K(s,\xi)$, определенная равенством $(\ref{K})$, неотрицательна.


\end{lems}





\begin{lems}


\label{l4}


Пусть оператор $T:C[a,b]\to L[a,b]$


имеет представление $(\ref{st})$ и изотонен. Пусть функция $K(s,\xi)$


определена равенством $(\ref{K})$. Тогда при любой ограниченной в


существенном неотрицательной функции $f$ функция


$$


\psi(s)\df


\int_a^bK(s,\xi)f(\xi)\,d\xi,\quad s\in[a,b],


$$


выпукла вверх.


\end{lems}





\begin{pf*}{Доказательство теоремы \ref{t7}}


Так как оператор $G$ антитонен, но не является строго


антитонным, то $G(t_0,\cdot)=0$ на


множестве ненулевой меры при некотором $t_0\in(a,b)$.





Из равенства (\ref{dzv}) и лемм \ref{l3}, \ref{l4} следует, что


если $G(t_0,s_0)=0$ при некоторых $t_0,s_0\in(a,b)$,


то либо $G(t_0,s)=0$ при


$s\in[a,s_0]$, либо $G(t_0,s)=0$ при


$s\in[s_0,b]$. Теперь из равенств (\ref{zv}) заключаем, что либо


$x_b(t_0)=0$ (при этом $x_b(t)=0$, если $t\in[a,t_0$)), либо


$y_a(t_0)=0$ (при этом $y_a(t)=0$, если $t\in[t_0,b$)). Здесь


$x_b$, $y_a$ --- решения задач (\ref{7}) и (\ref{8}) соответственно.


Отсюда следует утверждение теоремы.


\end{pf*}





\begin{notes}


\label{z6}


Пусть оператор $T:C[a,b]\to


L[a,b]$ имеет представление (\ref{st}), изотонен и выполнены


равенства (\ref{n9}). Предположим, что оператор Грина $G$ задачи


(\ref{11}) антитонен. Тогда,


если решение $x_b(\cdot)$ задачи (\ref{7})


таково, что $x_b(t)=0$ при $t\in[a,c]$ и $x_b(t)>0$ при $t\in(c,b]$


для некоторого $c\in(a,b)$, то


\begin{equation}\label{u1}


\text{


$r(t,s)=\const$ на отрезке $s\in[c,b]$ для почти всех


$t\in[a,c]$.}


\end{equation}


Если решение $y_a(\cdot)$ задачи (\ref{8}) таково, что


$y_a(t)=0$ при $t\in[c,b]$ и $y_a(t)>0$ при $t\in[a,c]$ для


некоторого $c\in(a,b)$, то


\begin{equation}\label{u2}


\text{$r(t,s)=\const$ на отрезке $s\in[a,c]$ для почти всех


$t\in[c,b]$.}


\end{equation}





Если вронскиан фундаментальной системы решений уравнения


$\L x=0$ имеет нуль в точке $c\in(a,b)$, то выполнено хотя бы


одно из двух условий (\ref{u1}), (\ref{u2}).


\end{notes}





Действительно, пусть для определенности


$x_b(t)=0$ при $t\in[a,c]$ и $x_b(t)>0$ при


$t\in(c,b]$ для некоторого $c\in(a,b)$. Так как $\pi\Ddot x_b=Tx_b$,


то $(Tx_b)(t)=0$ при $t\in[a,c]$. Так как $x_b(s)>0$ при


$s\in(c,b]$, то выполнено условие (\ref{u1}).





Остальные утверждения замечания \ref{z6} доказываются аналогично.


\medskip





Автор выражает свою признательность Н.В.\,Азбелеву и С.А.\,Гусаренко за их


внимание к работе и ценные замечания.





\begin{thebibliography} {10}





\bibitem{Lev} Левин А.Ю. {\em Неосцилляция решений уравнения


$x^{(n)}+p_1(t)x^{(n-1)}+\dotsb+p_n(t)=0$} //


Успехи матем. наук. -- 1969. -- Т.\,24. -- \No\,2. --


С.\,43--96.





\bibitem{KSh} Кигурадзе И.Т., Шехтер Б.Л. {\em Сингулярные краевые


задачи для обыкновенного дифференциального


уравнения второго порядка} // Итоги науки и техн.


Современ. пробл. матем. -- M.: ВИНИТИ, -- 1987. --


Т.\,30. -- С.\,105--201.





\bibitem{H} Хартман\,Ф. {\em Обыкновенные дифференциальные


уравнения}. -- М.: Мир, 1970. -- 720 с.





\bibitem{AD0} Азбелев Н.В., Домошницкий А.И. {\em О дифференциальном


неравенстве Валле--Пуссена} //


Дифференц. уравнения. -- 1986. -- Т.\,22.


-- \No\,12. -- С.\,2041--2045.





\bibitem{AMR} Азбелев Н.В., Максимов В.П., Рахматуллина


Л.Ф. {\em Введение в теорию функционально-дифференциальных уравнений}. --


М.: Наука, 1991. -- 280 с.





\bibitem{Lab} Лабовский С.М. {\em О положительных решениях


двухточечной краевой задачи для линейного


сингулярного функционально-дифференциального уравнения} //


Дифференц. уравнения. -- 1988. -- Т.\,24. --


\No\,10. -- С.\,1695--1704.





\bibitem{Kob} Кобяков И.И. {\em Условия отрицательности


функции Грина двухточечной краевой задачи с


отклоняющимся аргументом} // Дифференц. уравнения.


-- 1972. -- Т.\,8. -- \No\,3. -- С.\,443--452.





\bibitem{Pla} Плаксина В.П. {\em Условия знакопостоянства функции


Грина одной двухточечной задачи для функционально-дифференциального


уравнения $n$-го порядка}. -- Пермь, 1989. -- 43


с. Деп. в ВИНИТИ 16.05.89, \No\,3280-В89.





\bibitem{Lih} Лихачева Н.Н. {\em Теорема Валле--Пуссена для одного


класса функционально-дифференциальных


уравнений} // Изв. вузов. Математика. --


1997. -- \No\,5. -- С.\,30--37.





\bibitem{AD} Азбелев Н.В., Домошницкий А.И.


{\em К вопросу о линейных


дифференциальных неравенствах}. I //


Дифференц. уравнения. -- 1991. -- Т.\,27. -- \No\,3. --


С.\,376--384.





\bibitem{AD2} Азбелев Н.В., Домошницкий А.И.


{\em К вопросу о линейных


дифференциальных неравенствах}. II //


Дифференц. уравнения. -- 1991. -- Т.\,27. --


\No\,6. -- С.\,923--931.





\bibitem {AAB} Азбелев Н.В., Алвеш М., Бравый Е.И. {\em О сингулярных


краевых задачах для линейного


функционально-дифференциального уравнения


второго порядка} // Изв. вузов. Математика.


-- 1999. -- \No\,2. -- С.\,3--11.





\bibitem{Lab2} Лабовский С.М. {\em О сохранении знака


вронскиана фундаментальной системы, функции Коши


и функции Грина двухточечной краевой задачи для


уравнения с запаздывающим аргументом} //


Дифференц. уравнения. -- 1975.


-- Т.\,11. -- \No\,10. -- С.\,1780--1789.





\bibitem{KVP} Канторович Л.В., Вулих Б.З., Пинскер А.Г.


{\em Функциональный анализ в полуупорядоченных


пространствах}. -- М.--Л.: Гостехиздат, 1950. -- 548 c.





\bibitem{MR} Максимов В.П., Рахматуллина Л.Ф. {\em Линейное


функционально-дифференциальное уравнение, разрешенное относительно


производной} // Дифференц. уравнения. -- 1973. --


Т.\,9. -- \No\,12. -- С.\,2231--2240.





\bibitem{AR} Азбелев Н.В., Рахматуллина Л.Ф. {\em Об оценке


спектрального радиуса


линейного оператора в пространстве непрерывных


функций} // Изв. вузов. Математика. -- 1996.


-- \No\,11. -- С.\,23--28.





\end{thebibliography}





\vskip 0.5cm





\post{\it Пермский государственный}{\it Поступила}


\post{\it технический университет}{11.05.2000}





\label{end}


\end{document}





